\begin{textsample}{NEG Dim 6 – global\_south\_2025\_03 – Score -1.00 – global\_south\_2025\_03/122540905.txt}  \label{ex:f6_neg_140}
Khalil, a Columbia University graduate student and legal permanent resident, has been held by the government since March 8 in a push to deport him over his participation in campus protests for Gaza last year.
Student negotiator Mahmoud Khalil is on the Columbia University campus in New York at a pro-Palestinian protest encampment. ( AP ) A federal court in the United States has dismissed an effort by President Donald Trump ' s administration to dismiss Palestinian rights activist Mahmoud Khalil ' s legal challenge against his detention and deportation, but transferred the case to a court in New Jersey.
Khalil, a Columbia University graduate student and legal permanent resident, has been held by the government since March 8 in a push to deport him over his participation in campus protests for Gaza last year.
Manhattan-based US District Judge Jesse Furman denied the Trump administration ' s request to dismiss Khalil ' s case but determined that he lacked jurisdiction, as Khalil was being held in New Jersey when his lawyers first contested his arrest in New York. @ @ @ @ @ @ @ @ @ @ and Customs Enforcement ( ICE ) agents on March 8 outside his Columbia University \textbf{residence}.
His lawyers argue that his detention was politically motivated and violated his First Amendment rights.
In a letter from detention, Khalil referred to himself as a"political prisoner."He was initially held in Manhattan before being transferred to New Jersey, where his lawyer filed a habeas corpus petition.
However, Furman ruled that because Khalil was in New Jersey at the time, the case must be heard there."At the heart of this case is the important question of whether and under what circumstances the Government may rescind a person ' s lawful permanent resident status and remove him from the United States,"Furman wrote.
Story continues below this ad Khalil, a Palestinian descent student who entered the US on a visa in 2022 and obtained a green card last year, became a prominent figure in Columbia ' s pro-Palestinian movement.
The administration has sought his removal under a rarely used provision of the 1952 Immigration @ @ @ @ @ @ @ @ @ @ of State believes an individual ' s presence could harm US foreign policy.
Secretary of State Marco Rubio stated on March 16 that participating in"pro-Hamas events"contradicts US foreign policy.
Khalil ' s lawyers insist he has no ties to Hamas and was merely a mediator during the protests.
His wife, Noor Abdallah, an American citizen, is eight months pregnant and unable to visit him in Louisiana, where he is currently detained.
His lawyers argue that the government moved him there to prevent his case from being heard in New York or New Jersey.

% matched lemmas: residence
\end{textsample}
